\chapter{Calibración Geométrica}
\label{ch:calibration}
% ==============================================================

\section{¿Qué es la calibración y para qué sirve?}

Cada cámara ve la playa en \textbf{perspectiva}: un píxel de la imagen no dice, por sí solo,
qué punto real de la playa representa. La calibración resuelve exactamente ese problema: encuentra
la transformación matemática que convierte un píxel $(u,v)$ de la foto en una coordenada real
$(X,Y)$ en metros, en el sistema de referencia EPSG:25830 (ETRS89 / UTM huso 30N).

Esa transformación es una \textbf{homografía} $\mathbf{H}$, una matriz $3\times3$:

\begin{equation}
  \begin{pmatrix} X \\ Y \\ 1 \end{pmatrix}
  \sim
  \mathbf{H}
  \begin{pmatrix} u \\ v \\ 1 \end{pmatrix}
  \label{eq:homography}
\end{equation}

\begin{infobox}[¿Por qué una homografía y no algo más simple?]
Una homografía es la transformación correcta cuando se proyecta un plano (la arena de la
playa, que se asume aproximadamente plana) visto desde un punto de vista arbitrario —
mezcla en una sola matriz la rotación, traslación, escala \emph{y} la perspectiva de la
cámara. Un simple factor de escala metros/píxel no sirve aquí: un píxel cerca de la cámara
cubre muchísimo menos terreno real que un píxel cerca del horizonte, y esa relación cambia de
forma no lineal con la distancia — es exactamente lo que la homografía sí modela.
\end{infobox}

\begin{warnbox}[Ojo: esta NO es la única homografía del sistema]
Este capítulo cubre la homografía de \textbf{calibración geométrica}: píxel de la foto
$\rightarrow$ coordenada UTM real. Existe una \textbf{segunda homografía}, completamente
distinta, en el módulo de alineación masiva (capítulo~\ref{ch:batch}): esa es
\emph{imagen$\rightarrow$imagen} (corrige el pequeño desplazamiento de la cámara entre dos
fotos) y no tiene nada que ver con coordenadas UTM. Las dos se calculan con la misma función
de OpenCV (\texttt{cv2.findHomography} con RANSAC) pero resuelven problemas distintos, sobre
puntos de entrada distintos, y no se deben confundir. Ver la nota comparativa en
\S\ref{sec:dos-homografias}.
\end{warnbox}

\section{De dónde salen los puntos de control (GCPs)}

La homografía se ajusta a partir de \textbf{GCPs} (\emph{Ground Control Points}, aquí llamados
``varillas''): puntos de los que se conocen \emph{las dos} coordenadas del mismo punto físico —
dónde cae en la foto (píxel) y dónde está en la realidad (UTM, en metros). Hacen falta un
mínimo de 4 pares para que el sistema tenga suficiente información (una homografía tiene 8
grados de libertad); en la práctica se recomiendan 8 o más, y cuantos más se usen, más robusto
es el ajuste frente a un punto puntual mal marcado.

Cada GCP llega al sistema por una de estas dos vías:

\begin{itemize}[leftmargin=1.5em]
  \item \textbf{Manual}: clic directo sobre la imagen en la interfaz web (paso \emph{Marcación},
        \S\ref{sec:flujo-web}), con una lupa de precisión ($\times6$) para acertar el píxel exacto
        en fotos de varios miles de píxeles de ancho.
  \item \textbf{Importado de un levantamiento de campo} (CSV con columna \texttt{FILE} que liga
        cada varilla a la sesión de disparo en la que se topografió): el píxel llega como
        \emph{aproximado} (varilla ``pendiente'', marcador ámbar) y hay que confirmarlo con la
        lupa antes de que cuente para el ajuste — así un píxel mal ubicado en el CSV no puede
        colar error silenciosamente en la calibración.
\end{itemize}

\begin{infobox}[¿El cálculo se hace sobre una imagen o sobre todas las de la cámara?]
Sobre \textbf{una sola imagen}, la que esté seleccionada en ese momento en el paso Marcación
—cada imagen tiene su propio conjunto de varillas marcadas y su propia homografía calculada,
guardada junto a esa imagen concreta (\texttt{data/CAM\_X/json/\{imagen\}.json}, campo
\texttt{calibration})—. Lo que SÍ es por cámara es el \emph{resultado activo}: cada vez que se
calcula (o recalcula) la homografía de una imagen, esa $\mathbf{H}$ se copia también al
\textbf{perfil de la cámara} (\texttt{calibration/cam\_X\_profile.json} y
\texttt{cam\_X\_H.npy|}), sobrescribiendo la anterior. Es esa copia a nivel de cámara —la
última calibración calculada, sea de la imagen que sea— la que usa el análisis de línea de
costa (capítulo~\ref{ch:validation}) para \emph{todas} las imágenes nuevas de esa cámara, hasta
la siguiente recalibración. En otras palabras: se calibra imagen a imagen, pero solo hace falta
una calibración vigente por cámara para poder analizar el resto de sus capturas.
\end{infobox}

\section{Flujo actual: calibración interactiva desde la interfaz web}
\label{sec:flujo-web}

La calibración se hace hoy \textbf{enteramente desde el navegador} (pantalla
\emph{Calibración}), no con scripts de línea de comandos. El asistente tiene 7 pasos; los tres
que corresponden a este capítulo son \emph{Marcación}, \emph{Cálculo} y \emph{Validación}:

\begin{figure}[H]
\centering
\begin{tikzpicture}[node distance=0.5cm and 1.6cm]
  \node[block, fill=SkyBlue!15] (marc) {Marcación\\{\tiny marcar/confirmar varillas}};
  \node[block, fill=Orange!15, right=of marc] (calc) {Cálculo\\{\tiny RANSAC $\rightarrow \mathbf{H}$ + residuos}};
  \node[block, fill=OliveGreen!15, right=of calc] (val) {Validación\\{\tiny RMSE terreno vs.\ umbral}};

  \draw[arrow] (marc) -- (calc) node[midway,above,label]{$\geq 4$ confirmadas};
  \draw[arrow] (calc) -- (val) node[midway,above,label]{Continuar};
  \draw[arrow, bend left=35, color=BrickRed!60] (calc.north) to node[above,label]{\textcolor{BrickRed}{excluir + recalcular}} (calc.north);
\end{tikzpicture}
\caption{Los tres pasos del asistente de calibración implicados en este capítulo}
\label{fig:calib-steps}
\end{figure}

\begin{enumerate}[leftmargin=1.5em]
  \item \textbf{Marcación}: se marca (o confirma) cada varilla sobre la foto. Un clic directo
        encima del marcador de una varilla pendiente la confirma en su sitio; un clic en
        cualquier otro punto de la imagen la reposiciona. Cada varilla confirmada se pinta en
        \textbf{verde}; las pendientes, en ámbar.
  \item \textbf{Cálculo}: al pulsar \emph{Calcular homografía} se ejecuta el ajuste
        (\S\ref{sec:ransac}) y aparece una tabla de error de reproyección por varilla — se puede
        excluir una varilla concreta y recalcular sin ella, de forma iterativa, antes de pasar a
        Validación. Como la calibración es por imagen (\S\ref{sec:flujo-web}), este paso incluye
        también un selector para pinchar entre las distintas fotos de la cámara y comparar el
        resultado de cada una sin volver a Marcación; la última que se calcule es la que queda
        activa para la cámara. Cada fila de la tabla muestra, además del error, el
        \textbf{píxel y la coordenada UTM tal y como se marcaron} — para poder verificar cada
        varilla contra el CSV/Excel de origen si un valor no encaja como se esperaba — y, al
        pasar el ratón sobre el error, el valor \emph{predicho} por la homografía en ese mismo
        punto, para ver exactamente qué discrepancia produce el error.
  \item \textbf{Validación}: resumen final de la calidad de la calibración de esta imagen —
        RMSE en metros frente al umbral objetivo, y estado \emph{Óptimo}/\emph{Revisar}.
\end{enumerate}

\begin{warnbox}[Scripts de calibración retirados]
Versiones anteriores de este sistema incluían herramientas de línea de comandos
(\texttt{calibration\_tool.py}, \texttt{automate\_homography.py}, \texttt{recalibrate.py}) para
marcar GCPs fuera del navegador. Se han retirado del repositorio: la calibración de producción
vive por completo en el backend web (\texttt{calculate\_homography()}, invocado desde el paso
Cálculo) y esas herramientas habían quedado desincronizadas de ese flujo real —
\texttt{calibration\_tool.py} en concreto ni siquiera llegaba a compilar. El flujo de este
capítulo es el único vigente.
\end{warnbox}

\section{Estimación con RANSAC}
\label{sec:ransac}

\begin{codebox}[Estimación de homografía con RANSAC --- \texttt{calculate\_homography()} en el backend]
\begin{minted}{python}
H, ransac_mask = cv2.findHomography(
    pts_px,             # píxeles marcados/confirmados en la foto
    pts_utm,            # coordenadas UTM del catálogo de varillas
    cv2.RANSAC,
    3.0                 # umbral de reproyección en el dominio UTM: 3 metros
)
\end{minted}
\end{codebox}

\begin{infobox}[¿Qué hace RANSAC aquí y qué es un \emph{inlier}?]
RANSAC (\emph{Random Sample Consensus}) prueba subconjuntos de varillas y se queda con la
homografía que consigue que el mayor número posible de ellas encaje dentro de un margen de
error tolerable (aquí, 3~m en UTM). Las varillas que SÍ encajan con esa homografía se llaman
\textbf{inliers}; las que no —porque su UTM o su píxel están mal, o porque son incompatibles
con el resto— quedan como \emph{outliers} y no participan en el ajuste final. Es la misma idea
que un ``voto mayoritario'' geométrico: una sola varilla mal introducida no puede arruinar la
homografía si el resto de varillas son consistentes entre sí. La tabla de residuos del paso
Cálculo marca con un \checkmark verde las varillas que RANSAC aceptó como inlier y con una
\textcolor{BrickRed}{$\times$} roja las que descartó.
\end{infobox}

\subsection{Aviso de geometría inestable}
\label{sec:geometria-inestable}

Si las varillas marcadas quedan casi todas alineadas a lo largo de una única línea (por
ejemplo, una fila de estacas topografiadas siguiendo la orilla, con muy poca variación en la
dirección perpendicular), el sistema queda \textbf{numéricamente mal condicionado}: existen
infinitas homografías que explican igual de bien esos puntos, y una homografía así puede
fallar estrepitosamente (RMSE de decenas de metros) aunque el cálculo en sí no produzca ningún
error. Se detecta comprobando qué fracción de las varillas activas quedó como inlier de
RANSAC:

\begin{equation}
  f_\text{inlier} = \frac{\text{varillas inlier}}{\text{varillas activas}}
  \qquad\qquad
  f_\text{inlier} < 0{,}6 \;\Rightarrow\; \text{aviso de geometría inestable}
\end{equation}

Cuando esto ocurre, la interfaz muestra explícitamente cuántas varillas de cuántas encajaron y
recomienda añadir varillas repartidas en más de una zona/profundidad de la imagen, en vez de
dejar solo un RMSE alto sin explicación.

\section{Dos umbrales que no hay que confundir}
\label{sec:dos-umbrales}

La pantalla de calibración maneja \textbf{dos umbrales distintos}, en dominios distintos, con
propósitos distintos — es la fuente de confusión más habitual al leer los resultados:

\begin{table}[H]
\centering
\caption{Umbral por varilla (píxeles) frente a umbral de aceptación (metros)}
\label{tab:dos-umbrales}
\begin{tabularx}{\textwidth}{L{3.2cm} C{2.6cm} X}
\toprule
\textbf{} & \textbf{Umbral} & \textbf{Para qué sirve} \\
\midrule
Error por varilla & 5~px (ajustable) & Ayuda visual para \emph{localizar} qué varilla concreta
  podría estar mal puesta (un misclick, o una varilla del catálogo equivocada). No decide si la
  calibración es válida. \\
\addlinespace
RMSE terreno & 2~m (\texttt{RMSE\_GOOD\_M}) & \textbf{Criterio real de aceptación} de la
  calibración de esta imagen — el mismo valor que pondera el índice de confianza del análisis
  de línea de costa (capítulo~\ref{ch:validation}). \\
\bottomrule
\end{tabularx}
\end{table}

\begin{warnbox}[Por qué el estado ya NO depende de que todas las varillas bajen del umbral en píxeles]
En una iteración anterior de este sistema, el indicador ``Óptimo / Revisar'' exigía que
\emph{todas} las varillas individuales tuvieran un error de reproyección por debajo de 1~px.
Con fotos de varios miles de píxeles de ancho marcadas a mano, ese criterio era prácticamente
imposible de cumplir incluso con una calibración objetivamente buena: se observó un caso real
con RMSE~terreno~=~0,24~m (muy por debajo del objetivo de 2~m) donde 7 de 8 varillas aparecían
igualmente como ``sobre umbral''. El estado ahora se decide por el RMSE en metros frente a
\texttt{RMSE\_GOOD\_M}, que es la métrica que de verdad importa para georreferenciar la línea
de costa; el umbral en píxeles sigue existiendo solo como ayuda de diagnóstico por varilla.
\end{warnbox}

\section{Métricas del ajuste: RMSE y MAE}

Para cada varilla activa (no excluida, no pendiente) se calcula el error de reproyección en
los dos dominios:

\begin{align}
  e_{\text{px},i} &= \big\| \,\mathbf{H}^{-1}(\text{UTM}_i) - \text{píxel}_i \,\big\| \\
  e_{\text{m},i}  &= \big\| \,\mathbf{H}(\text{píxel}_i) - \text{UTM}_i \,\big\|
\end{align}

y a partir de esos residuos, el RMSE (penaliza más los errores grandes — útil para detectar
varillas mal marcadas) y el MAE (distancia media, sin elevar al cuadrado):

\begin{equation}
  \text{RMSE} = \sqrt{\frac{1}{N}\sum_{i=1}^N e_i^2}
  \qquad\qquad
  \text{MAE} = \frac{1}{N}\sum_{i=1}^N e_i
\end{equation}

\begin{warnbox}[Tabla histórica --- ver estado actual en \S\ref{sec:refinamiento-agosto31}]
Los resultados de esta sección corresponden a la campaña de calibración de junio 2026 y se
conservan como registro histórico. El estado y los RMSE actuales de las 6 cámaras, tras la
ronda de refinamiento del 31 de agosto de 2026 (incluida la corrección de cámara~6), están en
la tabla de \S\ref{sec:refinamiento-agosto31} (capítulo~\ref{ch:roadmap}).
\end{warnbox}

\section{Resultados de calibración (campaña junio 2026)}

\begin{table}[H]
\centering
\caption{Resultados de precisión por cámara — campaña de calibración inicial}
\label{tab:calibration-results}
\begin{tabular}{C{1.8cm} C{1.8cm} C{2.5cm} C{2.5cm} C{2.5cm} C{2.2cm}}
\toprule
\textbf{Cámara} & \textbf{GCPs} & \textbf{RMSE Cal. (m)} & \textbf{RMSE Val. (m)} & \textbf{RMSE Val. (px)} & \textbf{Estado} \\
\midrule
CAM~1 & 51 & 1.53 & \textbf{1.45} & 19.36 & \textcolor{OliveGreen}{OK} \\
CAM~2 & 39 & 2.64 & 3.45          & 21.09 & \textcolor{Goldenrod}{Revisar} \\
CAM~3 & 49 & 1.61 & 2.32          & 38.38 & \textcolor{OliveGreen}{OK} \\
CAM~4 & 49 & 2.78 & 3.06          & 100.01 & \textcolor{Goldenrod}{Revisar} \\
CAM~5 & 61 & 1.45 & \textbf{1.38} & 14.81 & \textcolor{OliveGreen}{OK} \\
CAM~6 & 69 & 6.14 & 4.74          & 130.62 & \textcolor{BrickRed}{Critico} \\
\bottomrule
\end{tabular}
\end{table}

\begin{warnbox}[Cámaras con error elevado]
Los errores de CAM~4 y CAM~6 superan ampliamente el objetivo de 2~m y son consistentes con el
problema descrito en \S\ref{sec:geometria-inestable}: en la sesión de campo correspondiente a
CAM~6, las varillas topografiadas quedan casi alineadas en una única línea a lo largo de la
orilla (variación de apenas $\sim$5~m en un eje frente a $\sim$140~m en el otro), lo que deja la
homografía numéricamente mal condicionada.
\textbf{Acción recomendada:} volver a topografiar varillas para estas cámaras repartidas en más
de una zona/profundidad de la imagen (no solo a lo largo de la orilla), y recalibrar desde el
paso Marcación de la interfaz web.
\end{warnbox}

\section{Imágenes de diagnóstico}

Las imágenes de la figura~\ref{fig:diagnostics} muestran la reproyección de GCPs sobre las
imágenes de referencia, de la campaña de calibración inicial. Los vectores de error se
representan magnificados $\times10$.

\begin{figure}[H]
  \centering
  \begin{subfigure}[b]{0.48\textwidth}
    \includegraphics[width=\textwidth]{../calibration/diagnostics/CAM1_diagnostic.jpg}
    \caption{CAM~1 — RMSE val. 1.45~m}
  \end{subfigure}
  \hfill
  \begin{subfigure}[b]{0.48\textwidth}
    \includegraphics[width=\textwidth]{../calibration/diagnostics/CAM2_diagnostic.jpg}
    \caption{CAM~2 — RMSE val. 3.45~m}
  \end{subfigure}

  \vspace{0.4cm}
  \begin{subfigure}[b]{0.48\textwidth}
    \includegraphics[width=\textwidth]{../calibration/diagnostics/CAM3_diagnostic.jpg}
    \caption{CAM~3 — RMSE val. 2.32~m}
  \end{subfigure}
  \hfill
  \begin{subfigure}[b]{0.48\textwidth}
    \includegraphics[width=\textwidth]{../calibration/diagnostics/CAM4_diagnostic.jpg}
    \caption{CAM~4 — RMSE val. 3.06~m \textcolor{Goldenrod}{(revisar)}}
  \end{subfigure}

  \vspace{0.4cm}
  \begin{subfigure}[b]{0.48\textwidth}
    \includegraphics[width=\textwidth]{../calibration/diagnostics/CAM5_diagnostic.jpg}
    \caption{CAM~5 — RMSE val. 1.38~m}
  \end{subfigure}
  \hfill
  \begin{subfigure}[b]{0.48\textwidth}
    \includegraphics[width=\textwidth]{../calibration/diagnostics/CAM6_diagnostic.jpg}
    \caption{CAM~6 — RMSE val. 4.74~m \textcolor{BrickRed}{(crítico)}}
  \end{subfigure}

  \caption{Diagnóstico de calibración — reproyección de GCPs (vectores $\times10$)}
  \label{fig:diagnostics}
\end{figure}

\section{Comparación con la homografía de alineación (capítulo~\ref{ch:batch})}
\label{sec:dos-homografias}

\begin{table}[H]
\centering
\caption{Las dos homografías del sistema, lado a lado}
\label{tab:dos-homografias}
\begin{tabularx}{\textwidth}{L{3.4cm} X X}
\toprule
\textbf{} & \textbf{Calibración (este capítulo)} & \textbf{Alineación masiva (cap.~\ref{ch:batch})} \\
\midrule
Dominio & Píxel de la foto $\rightarrow$ UTM (metros) & Píxel de una foto $\rightarrow$ píxel de
  otra foto (misma cámara) \\
Entrada & Varillas GCP marcadas a mano (píxel + UTM conocidos) & Puntos SIFT emparejados por
  FLANN entre dos capturas \\
Umbral RANSAC & 3~m (dominio UTM) & 5~px (dominio píxel) \\
Se calcula & Una vez por imagen de referencia, al calibrar & Una vez por cada imagen nueva del
  lote, contra la base \\
Para qué sirve & Georreferenciar la línea de costa detectada (obtener m$^2$, coordenadas reales) &
  Corregir el pequeño desplazamiento de la cámara (viento, mantenimiento) para que la
  calibración siga siendo válida en fotos posteriores \\
\bottomrule
\end{tabularx}
\end{table}

Ambas usan \texttt{cv2.findHomography(..., cv2.RANSAC, umbral)} porque el problema matemático
subyacente es el mismo (encontrar la transformación proyectiva que mejor explica un conjunto de
correspondencias de puntos, descartando las atípicas) — pero los puntos de entrada, el dominio
del umbral y el propósito final son distintos, y conviene no confundir un ``RMSE'' o unos
``inliers'' de una con los de la otra al leer los logs o la interfaz.

% ==============================================================
